\documentclass[10pt,a4paper]{amsart}
\usepackage{amsmath,amssymb,amsthm,graphicx}
% Retained preamble for the Stacks Project chapter snapshot in this directory.
% The chapter's own preamble.tex is a copy of the shared Stacks Project preamble; the block
% below is sliced MECHANICALLY, line for line, from it. Nothing was added, removed or
% redefined; the wrapper class, the xr-hyper/hyperref wiring and the externaldocument list
% are omitted because the assembler supplies the document class.
%
% lines taken from preamble.tex: 3, 5, 9, 11, 12, 156, 157, 158, ... (51 lines total)
%
% package coverage restored from the chapter's own preamble (2026-10-13): xy, xy-2cell

\usepackage[all]{xy}
\xyoption{2cell}
\UseAllTwocells

\newenvironment{reference}{\comment}{\endcomment}
\newenvironment{slogan}{\comment}{\endcomment}
\newenvironment{history}{\comment}{\endcomment}
\theoremstyle{plain}
\newtheorem{theorem}[subsection]{Theorem}
\newtheorem{proposition}[subsection]{Proposition}
\newtheorem{lemma}[subsection]{Lemma}
\theoremstyle{definition}
\newtheorem{definition}[subsection]{Definition}
\newtheorem{example}[subsection]{Example}
\newtheorem{exercise}[subsection]{Exercise}
\newtheorem{situation}[subsection]{Situation}
\theoremstyle{remark}
\newtheorem{remark}[subsection]{Remark}
\newtheorem{remarks}[subsection]{Remarks}
\numberwithin{equation}{subsection}
\def\lim{\mathop{\mathrm{lim}}\nolimits}
\def\colim{\mathop{\mathrm{colim}}\nolimits}
\def\Spec{\mathop{\mathrm{Spec}}}
\def\Hom{\mathop{\mathrm{Hom}}\nolimits}
\def\Ext{\mathop{\mathrm{Ext}}\nolimits}
\def\SheafHom{\mathop{\mathcal{H}\!\mathit{om}}\nolimits}
\def\SheafExt{\mathop{\mathcal{E}\!\mathit{xt}}\nolimits}
\def\Sch{\mathit{Sch}}
\def\Mor{\mathop{\mathrm{Mor}}\nolimits}
\def\Ob{\mathop{\mathrm{Ob}}\nolimits}
\def\Sh{\mathop{\mathit{Sh}}\nolimits}
\def\NL{\mathop{N\!L}\nolimits}
\def\CH{\mathop{\mathrm{CH}}\nolimits}
\def\proetale{{pro\text{-}\acute{e}tale}}
\def\etale{{\acute{e}tale}}
\def\QCoh{\mathit{QCoh}}
\def\Ker{\mathop{\mathrm{Ker}}}
\def\Im{\mathop{\mathrm{Im}}}
\def\Coker{\mathop{\mathrm{Coker}}}
\def\Coim{\mathop{\mathrm{Coim}}}
\DeclareMathSymbol{\boxtimes}{\mathbin}{AMSa}{"02}
\def\QCohstack{\mathcal{QC}\!\mathit{oh}}
\def\Cohstack{\mathcal{C}\!\mathit{oh}}
\def\Spacesstack{\mathcal{S}\!\mathit{paces}}
\def\Quotfunctor{\mathrm{Quot}}
\def\Hilbfunctor{\mathrm{Hilb}}
\def\Curvesstack{\mathcal{C}\!\mathit{urves}}
\def\Polarizedstack{\mathcal{P}\!\mathit{olarized}}
\def\Complexesstack{\mathcal{C}\!\mathit{omplexes}}
\def\Pic{\mathop{\mathrm{Pic}}\nolimits}
\def\Picardstack{\mathcal{P}\!\mathit{ic}}
\def\Picardfunctor{\mathrm{Pic}}
\def\Deformationcategory{\mathcal{D}\!\mathit{ef}}
\usepackage{cleveref}
\crefname{originallemma}{original lemma}{original lemmas}
\Crefname{originallemma}{Original lemma}{Original lemmas}
\makeatletter
\newlabel{lemma-base-change}{{lemma-base-change}{source}{}{sourcepointer-1}{}}
\newlabel{lemma-base-change@cref}{{[originallemma][0][]lemma-base-change}{[1][0][]source}{}{}{}}
\makeatother
\crefname{originallemma}{original lemma}{original lemmas}
\Crefname{originallemma}{Original lemma}{Original lemmas}
\makeatletter
\newlabel{lemma-pushforward-restriction}{{lemma-pushforward-restriction}{source}{}{sourcepointer-2}{}}
\newlabel{lemma-pushforward-restriction@cref}{{[originallemma][0][]lemma-pushforward-restriction}{[1][0][]source}{}{}{}}
\makeatother
\crefname{originallemma}{original lemma}{original lemmas}
\Crefname{originallemma}{Original lemma}{Original lemmas}
\makeatletter
\newlabel{sites-cohomology-section-base-change-map}{{0735}{source}{}{sourcepointer-3}{}}
\newlabel{sites-cohomology-section-base-change-map@cref}{{[originallemma][0][]0735}{[1][0][]source}{}{}{}}
\makeatother
\title{Stacks, Sheaves on Algebraic Stacks: lemma-base-change-higher-direct-images}
\author{The Stacks Project Authors}
\date{stacks-project:stacks-sheaves.tex:local-snapshot-20261012; GNU FDL 1.2 or later}
\begin{document}
\maketitle
\noindent\textit{Editorial note.} Statement span is lines 4146-4165 of sources/stacks-algebraic-stacks-sheaves/pinned-original-source.tex; proof span is lines 4167-4185, the complete author proof (17 body lines).\par\medskip
\noindent\textit{Editorial note.} Navigation stubs: 3 cited-but-not-included labels are supplied as non-load-bearing pointers with their pinned tags where the snapshot tag index has them.\par\medskip
\noindent\textit{Editorial note.} extra\_packages is empty on purpose: the retained preamble already loads xy with options and a second declaration clashes in the assembler wrapper.\par\medskip
\noindent\textit{Editorial note.} Text is retained verbatim, including the author's own typos and placeholders.\par\medskip
\noindent\textit{Editorial note.} Queue index 34 of queue\_algebraic-stacks-sheaves.json (pb=775, stmt 4146-4165, proof 4167-4185).\par\medskip
\noindent\textit{Allowed standard prerequisites (editorial).}\begin{itemize}
\item Explicit prerequisite boundary. The cited inputs are prerequisites; the argument in the retained span is not replaced by any cited result.
\end{itemize}\medskip
\ifdefined\hypertarget\hypertarget{sourcepointer-1}{}\fi
\noindent\textit{Source reference.} Original lemma lemma-base-change: sources/stacks-algebraic-stacks-sheaves/pinned-original-source.tex, source lines 765--780. Cited by the retained statement or proof but not part of the collected spans; supplied only as a navigation pointer. The argument does not depend on its content.\par\medskip
\ifdefined\hypertarget\hypertarget{sourcepointer-2}{}\fi
\noindent\textit{Source reference.} Original lemma lemma-pushforward-restriction: sources/stacks-algebraic-stacks-sheaves/pinned-original-source.tex, source lines 4092--4110. Cited by the retained statement or proof but not part of the collected spans; supplied only as a navigation pointer. The argument does not depend on its content.\par\medskip
\ifdefined\hypertarget\hypertarget{sourcepointer-3}{}\fi
\noindent\textit{Source reference.} Original lemma 0735: sources/stacks-algebraic-stacks-sheaves/pinned-original-source.tex, source lines 4185--4185. Cited by the retained statement or proof but not part of the collected spans; supplied only as a navigation pointer. The argument does not depend on its content.\par\medskip
\begin{lemma}
\label{lemma-base-change-higher-direct-images}
Let $S$ be a scheme. Let
$\tau \in \{Zariski,\linebreak[0] \etale,\linebreak[0]
smooth,\linebreak[0] syntomic,\linebreak[0] fppf\}$. Let
$$
\xymatrix{
\mathcal{Y}' \times_\mathcal{Y} \mathcal{X} \ar[r]_{g'} \ar[d]_{f'} &
\mathcal{X} \ar[d]^f \\
\mathcal{Y}' \ar[r]^g & \mathcal{Y}
}
$$
be a $2$-cartesian diagram of algebraic stacks over $S$. Then the base change
map is an isomorphism
$$
g^{-1}Rf_*\mathcal{F} \longrightarrow Rf'_*(g')^{-1}\mathcal{F}
$$
functorial for $\mathcal{F}$ in $\textit{Ab}(\mathcal{X}_\tau)$
or $\mathcal{F}$ in $\textit{Mod}(\mathcal{X}_\tau, \mathcal{O}_\mathcal{X})$.
\end{lemma}

\begin{proof}
The isomorphism $g^{-1}f_*\mathcal{F} = f'_*(g')^{-1}\mathcal{F}$ is
Lemma \ref{lemma-base-change} (and it holds for arbitrary presheaves).
For the total direct images, there is a base change map because the
morphisms $g$ and $g'$ are flat, see
Cohomology on Sites, Section \ref{sites-cohomology-section-base-change-map}.
To see that this map is a quasi-isomorphism we can use that for
an object $y'$ of $\mathcal{Y}'$ over a scheme $V$ there is an equivalence
$$
(\Sch/V)_{fppf} \times_{g(y'), \mathcal{Y}} \mathcal{X}
=
(\Sch/V)_{fppf} \times_{y', \mathcal{Y}'}
(\mathcal{Y}' \times_\mathcal{Y} \mathcal{X})
$$
We conclude that the induced map
$g^{-1}R^if_*\mathcal{F} \to R^if'_*(g')^{-1}\mathcal{F}$
is an isomorphism by
Lemma \ref{lemma-pushforward-restriction}.
\end{proof}

\end{document}
